\section{Google、Gemini 3.8 Live 系を発表}
Google公式はリアルタイム音声対話向け Gemini 3.8 Live と、複雑なタスク向け Extended Thinking を発表した。Search Live、Geminiアプリ、API公開プレビューなどへの展開を示し、97言語や非同期ツール呼び出しも紹介している。
\dailyxsource{Google AI (@GoogleAI)}{https://x.com/GoogleAI/status/2099908000924193124}

\section{ChatGPT / Codex から GPT-5.5 を退役予定}
OpenAIの製品アカウントは、10月14日にChatGPT、ChatGPT Work、Codexの全プランからGPT-5.5を外すと告知した。Codex利用者にはGPT-5.6 SolまたはGPT-6 Astraへの移行を案内する一方、API側では継続提供すると開発者アカウントが補足している。
\dailyxsource{ChatGPT (@ChatGPT)}{https://x.com/ChatGPT/status/2099954190600876533}

\section{OpenAIが今週とDevDayの大型出荷を予告}
Sam Altman とCodex / ChatGPT担当者は、今週の出荷とDevDayに向けて大きな発表があることを示唆した。観測窓内では具体的な製品名は示されておらず、予告としての話題である。
\dailyxsource{Sam Altman (@sama)}{https://x.com/sama/status/2099872600977760451}

\section{TypeSafe AI、意思決定特化モデル Jev を公開}
Diogo Almeida はTypeSafe AIのステルス解除とともに、文章生成ではなく型付き決定と確率を返す System One モデル Jev を公開した。20--200倍高速、40--400倍安価などの性能値は投稿者および公式ブログ側の主張で、独立検証は未確認である。
\dailyxsource{Diogo Almeida (@CompleteSkeptic)}{https://x.com/CompleteSkeptic/status/2099925682726002904}

\section{Salesforce、CRM推論モデル Koa を発表}
SalesforceはAgentforce向けCRM推論モデル Koa を発表した。NVIDIA Nemotron 3 Superを基盤とし、約27年分のCRM知見を模した合成データでポストトレーニングしたと説明する。顧客データは学習に使わず、まずパイロット提供するとしている。
\dailyxsource{Salesforce (@salesforce)}{https://x.com/salesforce/status/2099886699996160295}

\section{Gemini DeepThink 数学特化版とされるリーク}
コミュニティでは、Gemini DeepThinkの数学特化バリアントを開発中とする内部画面風の画像が拡散した。「Mathematica」と呼ぶ投稿もあるが、Googleからの公式確認はなく、画像の真正性も未検証である。
\dailyxsource{lyra (@lyraxana)}{https://x.com/lyraxana/status/2099916458549506315}

\section{Hugging Face エージェント事案の再解釈}
既存のHugging Face関連事案について、サンドボックス、評価設計、責任帰属、規制の妥当性をめぐる議論が続いた。観測窓内に新しい公式インシデント報告が出たわけではなく、数値やFTC調査への言及も投稿者主張として扱う。
\dailyxsource{Brian Roemmele (@BrianRoemmele)}{https://x.com/BrianRoemmele/status/2099871292619215200}

\section{RLの ``Matthew Effect'' と Never Giving Up}
Michael Noukhovitch は、LLMへのRL利得が易しい問題へ偏る ``Matthew Effect'' を報告し、難しい問題を解くための非同期RL手法 Never Giving Up を提案した。arXiv論文は提示されているが、独立再現は未確認である。
\dailyxsource{Michael Noukhovitch (@mnoukhov)}{https://x.com/mnoukhov/status/2099933311552397424}

\section{物語学習からアシスタント行動が転移するとの報告}
Owain Evans は、合成した人間ストーリーのみで学習したモデルが、通常チャットでも登場人物の特徴的な行動を示したと報告した。観測時点では論文PDFの確定URLが取得できておらず、著者による研究紹介として扱う。
\dailyxsource{Owain Evans (@OwainEvans\_UK)}{https://x.com/OwainEvans_UK/status/2099896330009391269}

\section{エージェント性能を Elo-per-token で分析}
Qiuyang Mang らは、長時間のテスト時計算に対するエージェント性能を Elo-per-token 曲線で分析したと報告した。長期では対数線形に収束し、単一の長時間ランより分割・並列実行が有利との主張だが、論文URLは観測時点で未確認である。
\dailyxsource{Qiuyang Mang (@MangQiuyang)}{https://x.com/MangQiuyang/status/2099917781529694431}

\section{Claude Code の更新を非公式アカウントが観測}
非公式観測アカウントは、Claude DesktopのSessions hubやClaude Code 2.1.273の変更点を報告した。リモートセッションのフォークやMCP切断通知などが挙げられているが、Anthropic公式チェンジログとの突合は未了である。
\dailyxsource{TestingCatalog (@testingcatalog)}{https://x.com/testingcatalog/status/2099933727778279803}

\section{Vercel v0 がモデル選択を拡大}
Vercelのv0公式は、AI Gateway経由でClaude、GPT、Kimi、GLM、Grok、DeepSeekなど複数モデルを選べるようになったと発表した。料金や既定モデルの変更など詳細は観測時点で未確認である。
\dailyxsource{v0 (@v0)}{https://x.com/v0/status/2099902412253499478}

\section{Grok Build CLI 1.0.33 のコミュニティ観測}
Grok追跡アカウントが、Grok Build CLI 1.0.33の変更としてMCP構造化JSON、メモリモーダル、複数の不具合修正を紹介した。xAI公式のリリースノートは本収集では確認されておらず、二次情報である。
\dailyxsource{Puck (@GrokInsider)}{https://x.com/GrokInsider/status/2099965928154812908}

\section{LeCun、LLM単体路線への懐疑を再主張}
Yann LeCun はLLMツールの有用性を認めつつ、人間レベルAIへの経路ではなく、実世界理解には異なるアーキテクチャが必要だと改めて主張した。新モデル発表ではなく、本人の研究上の立場表明である。
\dailyxsource{Yann LeCun (@ylecun)}{https://x.com/ylecun/status/2099806809431081258}
